\begin{knowbox}
\textbf{Nội dung:} Chương 1 -- 13 (tỉ số phần trăm, hình hộp chữ nhật và hình lập phương, thể tích, sơn bề mặt, chuyển động đều, thống kê và xác suất, sơ đồ Venn, tính ngược, giả thiết tạm). \quad \textbf{Thời gian gợi ý:} $60$ phút.
\end{knowbox}
\gap
\begin{probbox}{1}{\relax}
Một lớp học có $25$ bạn nữ và $15$ bạn nam. \pa{a} Số bạn nữ chiếm bao nhiêu phần trăm số học sinh cả lớp? \pa{b} Số bạn nam bằng bao nhiêu phần trăm số bạn nữ?
\end{probbox}
\probsol{
Cả lớp có $25+15=40$ học sinh.\par
\pa{a} $25:40=0{,}625=62{,}5\%$.\par
\pa{b} $15:25=0{,}6=60\%$.
\dapso{\pa{a} $62{,}5\%$;\ \pa{b} $60\%$.}}
\gap
\begin{probbox}{2}{\relax}
Cô Hạnh gửi tiết kiệm $5\,000\,000$ đồng với lãi suất $0{,}5\%$ một tháng. Sau một tháng cô Hạnh nhận được bao nhiêu tiền lãi, và khi đó tổng số tiền gửi và tiền lãi là bao nhiêu?
\end{probbox}
\probsol{
Tiền lãi: $5\,000\,000\times0{,}5:100=25\,000$ đồng.\par
Tổng số tiền: $5\,000\,000+25\,000=5\,025\,000$ đồng.
\dapso{$25\,000$ đồng; $5\,025\,000$ đồng.}}
\gap
\begin{probbox}{3}{\relax}
Ngày thứ nhất Khánh đọc được $40\%$ số trang của một quyển sách, ngày thứ hai đọc được $35\%$ số trang của quyển sách đó. Quyển sách còn $75$ trang chưa đọc. Hỏi quyển sách có bao nhiêu trang?
\end{probbox}
\probsol{
Số trang còn lại chiếm: $100\%-40\%-35\%=25\%$ số trang sách.\par
Quyển sách có: $75:25\times100=300$ trang.
\dapso{$300$ trang.}}
\gap
\begin{probbox}{4}{\stars{1}}
Có $200$ g nước muối chứa $5\%$ muối. Cần đổ thêm bao nhiêu gam nước lã để được loại nước muối chứa $4\%$ muối?
\end{probbox}
\probsol{
Lượng muối: $200\times5:100=10$ g; đổ thêm nước lã thì lượng muối không đổi.\par
Trong nước muối mới, $10$ g muối chiếm $4\%$, nên nước muối mới nặng $10:4\times100=250$ g.\par
Lượng nước cần đổ thêm: $250-200=50$ g.
\dapso{$50$ g.}}
\gap
\begin{probbox}{5}{\relax}
Một chiếc hộp không có nắp dạng hình hộp chữ nhật dài $30$ cm, rộng $20$ cm, cao $15$ cm. Tính diện tích cần sơn mặt ngoài của chiếc hộp.
\end{probbox}
\probsol{
Diện tích xung quanh: $(30+20)\times2\times15=1\,500$ cm$^2$.\par
Diện tích đáy: $30\times20=600$ cm$^2$.\par
Diện tích cần sơn: $1\,500+600=2\,100$ cm$^2$.
\dapso{$2\,100$ cm$^2$.}}
\gap
\begin{probbox}{6}{\stars{1}}
Một hình hộp chữ nhật có diện tích xung quanh là $200$ cm$^2$, chiều cao $5$ cm. Chiều dài hơn chiều rộng $4$ cm. Tính thể tích hình hộp chữ nhật đó.
\end{probbox}
\probsol{
Chu vi mặt đáy: $200:5=40$ cm, nên chiều dài $+$ chiều rộng $=20$ cm.\par
Chiều dài: $(20+4):2=12$ cm; chiều rộng: $12-4=8$ cm.\par
Thể tích: $12\times8\times5=480$ cm$^3$.
\dapso{$480$ cm$^3$.}}
\gap
\begin{probbox}{7}{\relax}
Một bể bơi dài $25$ m, rộng $10$ m, sâu $1{,}6$ m. Một máy bơm mỗi giờ bơm được $50$ m$^3$ nước. Bể đang cạn, hỏi phải bơm trong bao lâu thì mực nước bằng $80\%$ chiều sâu của bể?
\end{probbox}
\probsol{
Chiều cao mực nước: $1{,}6\times80\%=1{,}28$ m.\par
Thể tích nước: $25\times10\times1{,}28=320$ m$^3$.\par
Thời gian bơm: $320:50=6{,}4$ giờ $=6$ giờ $24$ phút.
\dapso{$6$ giờ $24$ phút.}}
\gap
\begin{probbox}{8}{\relax}
Một khối gỗ hình hộp chữ nhật kích thước $3$ cm $\times$ $4$ cm $\times$ $5$ cm được sơn tất cả các mặt, rồi cắt thành các hình lập phương nhỏ cạnh $1$ cm. Có bao nhiêu hình lập phương nhỏ được sơn \pa{a} $3$ mặt; \pa{b} $2$ mặt; \pa{c} $1$ mặt; \pa{d} không được sơn mặt nào?
\end{probbox}
\probsol{
Có tất cả $3\times4\times5=60$ hình lập phương nhỏ.\par
\pa{a} $8$ hình ở $8$ đỉnh.\par
\pa{b} Các hình ở giữa các cạnh: mỗi kích thước có $4$ cạnh, nên có $4\times(1+2+3)=24$ hình.\par
\pa{c} Phần trong của các mặt: $2\times(1\times2+2\times3+1\times3)=2\times11=22$ hình.\par
\pa{d} Phần lõi: $1\times2\times3=6$ hình. \ (\emph{Thử lại:} $8+24+22+6=60$.)
\dapso{$8$; $24$; $22$; $6$.}}
\gap
\begin{probbox}{9}{\relax}
Cần bao nhiêu hình lập phương nhỏ cạnh $1$ cm để xếp thành một khối lập phương đặc cạnh $5$ cm? Nếu khối lập phương đó rỗng bên trong, chỉ xếp lớp vỏ ngoài dày một hình lập phương nhỏ, thì cần bao nhiêu hình?
\end{probbox}
\probsol{
Khối đặc: $5\times5\times5=125$ hình lập phương nhỏ.\par
Phần rỗng bên trong là khối lập phương cạnh $5-2=3$ cm: $3\times3\times3=27$ hình.\par
Lớp vỏ ngoài: $125-27=98$ hình lập phương nhỏ.
\dapso{$125$ hình; $98$ hình.}}
\gap
\begin{probbox}{10}{\relax}
Lan đi bộ từ nhà lúc $6$ giờ $45$ phút đến nhà ông bà cách đó $4{,}5$ km với vận tốc $5$ km/giờ. Hỏi Lan đến nhà ông bà lúc mấy giờ?
\end{probbox}
\probsol{
Thời gian đi: $4{,}5:5=0{,}9$ giờ $=54$ phút.\par
Lan đến lúc: $6$ giờ $45$ phút $+\ 54$ phút $=7$ giờ $39$ phút.
\dapso{Lúc $7$ giờ $39$ phút.}}
\gap
\begin{probbox}{11}{\stars{1}}
Hai người ở cách nhau $27$ km đi bộ ngược chiều về phía nhau với vận tốc $4$ km/giờ và $5$ km/giờ. Một chú chó xuất phát cùng người thứ nhất, chạy qua chạy lại giữa hai người với vận tốc $10$ km/giờ cho đến khi hai người gặp nhau. Hỏi chú chó đã chạy được bao nhiêu ki-lô-mét?
\end{probbox}
\probsol{
Thời gian để hai người gặp nhau: $27:(4+5)=3$ giờ.\par
Chú chó chạy suốt thời gian đó: $10\times3=30$ km.
\dapso{$30$ km.}}
\gap
\begin{probbox}{12}{\relax}
Ô tô A ở phía trước ô tô B $60$ km, hai xe đi cùng chiều. Ô tô A đi với vận tốc $40$ km/giờ, ô tô B đi với vận tốc $55$ km/giờ. Hỏi sau bao nhiêu giờ ô tô B đuổi kịp ô tô A?
\end{probbox}
\probsol{
Mỗi giờ ô tô B gần ô tô A thêm $55-40=15$ km.\par
Thời gian đuổi kịp: $60:15=4$ giờ.
\dapso{$4$ giờ.}}
\gap
\begin{probbox}{13}{\stars{1}}
Một xe tải đi quãng đường $120$ km: $60$ km đầu với vận tốc $30$ km/giờ, $60$ km sau với vận tốc $60$ km/giờ. Tính vận tốc trung bình của xe tải trên cả quãng đường.
\end{probbox}
\probsol{
Thời gian đi: $60:30+60:60=2+1=3$ giờ.\par
Vận tốc trung bình: $120:3=40$ km/giờ (không phải $(30+60):2=45$).
\dapso{$40$ km/giờ.}}
\gap
\begin{probbox}{14}{\relax}
Một ca nô có vận tốc khi nước lặng là $14$ km/giờ, vận tốc dòng nước là $2$ km/giờ. Bến B ở phía hạ lưu, cách bến A $48$ km. Hỏi ca nô đi từ A đến B rồi quay về A hết bao nhiêu thời gian?
\end{probbox}
\probsol{
Xuôi dòng: $14+2=16$ km/giờ, hết $48:16=3$ giờ.\par
Ngược dòng: $14-2=12$ km/giờ, hết $48:12=4$ giờ.\par
Tổng thời gian: $3+4=7$ giờ.
\dapso{$7$ giờ.}}
\gap
\begin{probbox}{15}{\relax}
Hai đoàn tàu chạy ngược chiều trên hai đường ray song song: đoàn tàu thứ nhất dài $150$ m, vận tốc $20$ m/giây; đoàn tàu thứ hai dài $130$ m, vận tốc $15$ m/giây. Tính thời gian từ lúc hai đầu tàu gặp nhau đến lúc hai đuôi tàu rời nhau.
\end{probbox}
\probsol{
Hai tàu phải đi tổng quãng đường bằng tổng chiều dài: $150+130=280$ m.\par
Mỗi giây hai tàu đi được $20+15=35$ m.\par
Thời gian: $280:35=8$ giây.
\dapso{$8$ giây.}}
\gap
\begin{probbox}{16}{\stars{1}}
Trên một đường chạy hình tròn, Hòa chạy một vòng hết $6$ phút, Dũng chạy một vòng hết $10$ phút. Hai bạn xuất phát cùng lúc, cùng một điểm và chạy cùng chiều. Hỏi sau bao nhiêu phút Hòa đuổi kịp Dũng lần đầu? Khi đó Hòa đã chạy được bao nhiêu vòng?
\end{probbox}
\probsol{
Mỗi phút Hòa chạy được $\dfrac16$ vòng, Dũng chạy được $\dfrac1{10}$ vòng, nên Hòa hơn Dũng $\dfrac16-\dfrac1{10}=\dfrac1{15}$ vòng.\par
Để hơn đúng một vòng: $1:\dfrac1{15}=15$ phút.\par
Số vòng Hòa chạy: $15:6=2{,}5$ vòng.
\dapso{$15$ phút; $2{,}5$ vòng.}}
\gap
\begin{probbox}{17}{\relax}
\pa{a} Trong một bài kiểm tra có $3$ bạn được $6$ điểm, $5$ bạn được $7$ điểm, $8$ bạn được $8$ điểm và $4$ bạn được $9$ điểm. Tính điểm trung bình.\par
\pa{b} Gieo một con xúc xắc $50$ lần, có $8$ lần xuất hiện mặt $6$ chấm. Số lần xuất hiện mặt $6$ chấm chiếm bao nhiêu phần trăm tổng số lần gieo?
\end{probbox}
\probsol{
\pa{a} Số học sinh: $3+5+8+4=20$.\par Tổng điểm: $3\times6+5\times7+8\times8+4\times9=18+35+64+36=153$.\par Điểm trung bình: $153:20=7{,}65$.\par
\pa{b} $8:50=0{,}16=16\%$.
\dapso{\pa{a} $7{,}65$;\ \pa{b} $16\%$.}}
\gap
\begin{probbox}{18}{\relax}
Trong một nhóm $50$ học sinh có $30$ bạn học tiếng Anh, $26$ bạn học tin học và $10$ bạn không học môn nào trong hai môn đó. Hỏi có bao nhiêu bạn học cả hai môn?
\end{probbox}
\probsol{
Số bạn học ít nhất một môn: $50-10=40$.\par
Số bạn học cả hai môn: $30+26-40=16$.
\dapso{$16$ bạn.}}
\gap
\begin{probbox}{19}{\stars{1}}
Ngày thứ nhất, một gia đình dùng một nửa lượng nước trong bể và thêm $20$ lít. Ngày thứ hai, gia đình dùng một nửa lượng nước còn lại và thêm $10$ lít. Trong bể còn $50$ lít. Hỏi lúc đầu bể có bao nhiêu lít nước?
\end{probbox}
\probsol{
Tính ngược từ cuối lên.\par
Trước ngày thứ hai: $(50+10)\times2=120$ lít.\par
Lúc đầu: $(120+20)\times2=280$ lít.
\dapso{$280$ lít.}}
\gap
\begin{probbox}{20}{\stars{1}}
Một bài kiểm tra có $30$ câu: mỗi câu đúng được $4$ điểm, mỗi câu sai bị trừ $1$ điểm. Sơn trả lời tất cả các câu và được $95$ điểm. Hỏi Sơn trả lời đúng bao nhiêu câu?
\end{probbox}
\probsol{
Nếu đúng hết: $4\times30=120$ điểm; Sơn kém đi $25$ điểm.\par
Mỗi câu sai làm mất $4+1=5$ điểm: số câu sai $25:5=5$; số câu đúng $30-5=25$.
\dapso{$25$ câu.}}
